%======================================================================================
\section{Graphs, from scratch}\label{app:graphs}

\begin{itemize}
  \item \textbf{Graph.} A set of \emph{nodes} and a set of \emph{edges}, each edge joining two
        nodes. Here, at each instant, a node is a flight and an edge a pair interval in progress.
  \item \textbf{Path.} A sequence of nodes, each joined by an edge to the next. Two nodes are
        \emph{connected} when a path joins them.
  \item \textbf{Connected component.} A largest set of nodes connected to each other. Every node is
        in exactly one component; a node with no edge is a component alone (a pilot flying solo).
        A \emph{configuration} is a component with at least two nodes.
\end{itemize}

Two members of a configuration need not share an edge: in the chain 1--2--3 of
Figure~\ref{fig:graph}, pilots 1 and 3 are in the same configuration through pilot 2.

\begin{figure}[H]
\centering
\begin{tikzpicture}[x=1cm,y=1cm, node/.style={circle, draw=ink, fill=white, minimum size=6mm, inner sep=0pt,
    font=\footnotesize}, edge/.style={ink, line width=0.9pt}]
  \fill[blueI!10] (-0.5,-0.5) rectangle (2.5,1.5);
  \fill[blueI!10] (3.2,-0.5) rectangle (4.8,1.5);
  \node[node] (a) at (0,0) {1}; \node[node] (b) at (1,1) {2}; \node[node] (c) at (2,0) {3};
  \node[node] (d) at (3.5,0) {4}; \node[node] (e) at (4.5,1) {5}; \node[node] (f) at (6,0.5) {6};
  \draw[edge] (a) -- (b) -- (c); \draw[edge] (d) -- (e);
  \node[font=\footnotesize, text=ink] at (1,-0.85) {configuration $\{1,2,3\}$ (a chain)};
  \node[font=\footnotesize, text=ink] at (4,-0.85) {$\{4,5\}$};
  \node[font=\footnotesize, text=ink] at (6,-0.85) {solo};
\end{tikzpicture}
\caption{The graph at one instant: three components, two of them configurations.}
\label{fig:graph}
\end{figure}

\begin{itemize}
  \item \textbf{Directed graph.} The edges carry an arrow. For the episodes the nodes are
        configurations, with an arrow $c\to c'$ when $c'$ starts as $c$ ends and they share a member.
        An episode is a connected component of this graph, arrows ignored.
  \item \textbf{Bipartite graph.} Nodes of two kinds, edges only between kinds. At one instant, the
        configurations ending and those starting, with an edge when they share a member. Each
        component is one event (Figure~\ref{fig:events}); its outcome follows from how many
        configurations it has on each side: one to one is a join or a leave (members added or
        removed), one to several a split, several to one a merge.
\end{itemize}

\begin{figure}[H]
\centering
\begin{tikzpicture}[x=1cm,y=1cm, conf/.style={rounded corners=2pt, fill=blueI, text=white, font=\footnotesize,
    minimum height=5mm, inner xsep=4pt}, link/.style={->, ink, line width=0.8pt}]
  \node[conf] (c1) at (0,0.6) {$\{1,2\}$}; \node[conf] (c2) at (0,-0.6) {$\{3,4\}$};
  \node[conf] (c3) at (2.6,0) {$\{1,2,3,4\}$};
  \draw[link] (c1.east) -- (c3.west); \draw[link] (c2.east) -- (c3.west);
  \node[font=\footnotesize, text=ink] at (1.3,-1.3) {merge: two to one};
  \node[conf] (d1) at (6,0) {$\{1,2,3\}$};
  \node[conf] (d2) at (8.4,0) {$\{1,2\}$};
  \draw[link] (d1.east) -- (d2.west);
  \node[font=\footnotesize, text=ink] at (7.2,-1.3) {one to one, 3 removed: leave};
  \node[font=\footnotesize, text=ink] at (0,1.3) {ending}; \node[font=\footnotesize, text=ink] at (2.6,1.3) {starting};
\end{tikzpicture}
\caption{Two events: each is a connected component of the bipartite graph of the configurations
ending and starting at one instant.}
\label{fig:events}
\end{figure}

\textbf{Shape} (Section~\ref{sec:chains}). For a connected graph of $n$ nodes and $m$ edges:
\begin{itemize}
  \item \textbf{Density} $m/\binom{n}{2}$: the share of the node pairs joined by an edge.
  \item \textbf{Distance} between two nodes: the fewest edges of a path joining them. \textbf{Diameter}: the
        largest distance. In the chain 1--2--3 the diameter is 2; in a triangle, 1.
  \item \textbf{Complete graph}: every pair joined, density 1, diameter 1. \textbf{Tree}: $m=n-1$, the fewest
        edges that keep $n$ nodes connected, with no closed path; a \emph{chain} is a tree with diameter
        $n-1$, a \emph{star} one with diameter 2.
\end{itemize}

\textbf{Order} (Section~\ref{sec:leaders}). In a directed graph:
\begin{itemize}
  \item a \textbf{cycle} is a path that follows the arrows back to its start (A$\to$B$\to$C$\to$A);
        a graph with none is \textbf{acyclic}, and its nodes can be put in a line with every arrow
        pointing forward;
  \item a \textbf{source} has no arrow coming in: here, a pilot with nobody ahead, a leader;
  \item the \textbf{depth} of an acyclic graph is its longest chain of arrows: 1 for a star, $n-1$ for a line.
\end{itemize}
A directed triangle is \emph{transitive} (A$\to$B, B$\to$C, A$\to$C) or \emph{cyclic}; with arrows drawn
at random, 2 of its 8 orientations are cyclic.

\textbf{Configuration model} (Section~\ref{sec:partners}). A random graph with the \emph{degrees} (the number
of edges of each node) of a given one: each node gets as many half-edges as its degree, and the
half-edges are paired at random~\cite{newman2010}. It keeps how many partners each pilot had and
changes only who they were.

%======================================================================================
\section{Two models}\label{app:models}

\textbf{Ornstein--Uhlenbeck process}~\cite{uhlenbeck1930} (Section~\ref{sec:cohesion}). A quantity $\ell$
pulled back to 0 by a spring of rate $k$ and shaken by noise:
\begin{equation}
  d\ell=-k\,\ell\,dt+\sqrt{2k}\,\sigma\,dW .
\end{equation}
It forgets its value after $1/k$ and spreads by $\sigma$ around 0. Its mean square change after a lag
$s$ is $D(s)=\langle(\ell(t+s)-\ell(t))^2\rangle=2\sigma^2(1-e^{-ks})$: it grows as $2k\sigma^2 s$ at short
lags, like a free random walk, and stops at $2\sigma^2$ after $1/k$. With $k=0$ it is a random walk and
$D$ grows forever. A measurement noise of variance $c/2$ on each value adds $c$ to $D(s)$ at every lag.

\textbf{Birth--death chain} (Section~\ref{sec:fission}). A group of $n$ pilots gains one at rate $g(n)$ and
loses one at rate $l(n)$. Over a long time the transitions from $n$ to $n+1$ and back are equally
many, so the time spent at $n+1$ and at $n$, $T(n+1)$ and $T(n)$, satisfy $T(n)\,g(n)=T(n+1)\,l(n+1)$,
Eq.~\eqref{eq:chain}. Multiplying the ratios from 2 to $n$ gives $T(n)$. If every member joins others at
rate $a$ and leaves at rate $\mu$, $g(n)=an$, $l(n)=\mu n$ and $T(n)\propto x^n/n$ with $x=a/\mu$, the
log-series: the tail falls exponentially. A leaving rate that falls with $n$ makes it heavier.

%======================================================================================
\section{Phase at entry and exit}\label{app:phases}

\begin{table}[H]
\centering\footnotesize
\begin{tabular}{@{}l l l l@{}}
\toprule
 & climb (\%) & search (\%) & transition (\%) \\
\midrule
entering a group & 28.1 [27.6, 28.5] & 10.3 [10.0, 10.6] & 61.5 [61.0, 61.9] \\
leaving a group & 28.0 [27.6, 28.3] & 11.2 [10.9, 11.5] & 60.6 [60.1, 61.0] \\
share of time, alone & 29.1 [28.9, 29.3] & 10.1 [10.0, 10.2] & 60.8 [60.6, 61.0] \\
share of time, in a group & 33.0 [32.8, 33.2] & 9.4 [9.2, 9.6] & 57.6 [57.3, 57.9] \\
\bottomrule
\end{tabular}

\caption{Phase at the instant of entering or leaving a group, and shares of time in each phase
(percent, with 5--95\% bands; 1~Hz flights).}
\label{tab:phases}
\end{table}

%======================================================================================
\section{Reproducing the report}\label{app:reproduce}

From \texttt{analysis/groups/group-identification}, after the scripts of \emph{Group and solo flights}:
\begin{verbatim}
uv run python describe.py     # sizes, durations, events; about 15 s
uv run python cohesion.py     # delays and separations, one pass over every date; about 25 min
uv run python chains.py       # shapes and their H_loc; about 40 s
uv run python hierarchy.py    # order, switches, who leads; after cohesion.py
uv run python partners.py     # repeated partners; about 4 min
uv run python fission.py      # the birth-death chain; about 10 s
uv run python report/tables.py; uv run python report/structure_tables.py
cd report && latexmk -pdf group_structure.tex
\end{verbatim}
The per-date cache of \texttt{cohesion.py} is in \texttt{analyses/cohesion/all/days/} of the run on SSD\_DISANTE.
